% !TEX root = ../main.tex
\section{ROCKET Method}
\label{sec:rocket}
ROCKET transforms time series by applying convolutional kernels, which are not learned through training, but instead are randomly initialized \cite{dempster2020rocket}. Each kernel has parameters including length, weights, dilation, and bias, all of which are assigned randomly. The randomness allows the algorithm to generate a diverse set of filters, capable of capturing a wide variety of temporal patterns.

Once the kernels are applied to a time series, each convolution produces a feature map, showing the activation of the kernel across the entire series. From each feature map, ROCKET extracts two key summary features: Maximum value and Proportion of Positive Values (PPV).

Using a large number of kernels, ROCKET transforms each time series into a high-dimensional feature vector, which captures a large variety of temporal patterns without requiring kernel training.

These feature vectors are then fed into a linear classifier, typically Ridge Regression or Logistic Regression, which performs the final classification.

\subsection{Kernels}
ROCKET relies on a large number of convolutional kernels to extract meaningful features from time series. Each kernel acts like a small filter that highlights certain patterns in the data. Although the kernels are randomly initialized and not learned, their diversity allows the algorithm to extract rich and varied features from the time series. This random initialization enables ROCKET to efficiently explore a wide variety of temporal patterns without the computational cost of training each kernel.

The main parameters of each kernel are:
\begin{itemize}
    \item Length $l$: determines how many consecutive points the kernel considers, which allows kernels of different scales to detect both short and long-term patterns.
    \item Weights $w$: sampled from a standard normal distribution, giving each kernel a unique pattern of emphasis over the input series.
    \item Dilation $d$: controls the spacing between points in the kernel. A dilation of $d > 1$ allows the kernel to skip points in the series, allowing the model to capture patterns that occur over longer time intervals or at multiple scales.
    \item Bias $b$: random offset added to the convolution output, which increases variability of the resulting feature map.
\end{itemize}

Additionally, padding is applied when necessary to ensure that kernels cover the full series, preserving feature extraction across all time points. Each kernel produces a feature map along the series, which will later be summarized in the feature extraction stage.

\subsection{Time Series Transformation}
The core of the ROCKET algorithm lies in the transformation of the input time series via the application of the randomly generated convolutional kernels. Given a univariate time series $X = (x_1, x_2, \dots, x_n)$ and a random kernel with weights $W = (w_1, w_2, \dots, w_l)$, length $l$, dilation $d$, and bias $b$, the time series transformation is achieved through a 1D convolution operation.The kernel slides across the time series, computing the dot product between its weights and the corresponding subsequences of the input. Because of the dilation parameter $d$, the kernel does not necessarily interact with strictly contiguous data points; rather, it spans a receptive field of $(l - 1)d + 1$. For a specific position $i$ in the time series, the convolution output (activation) is calculated as:$$a_i = \left( \sum_{j=0}^{l-1} x_{i+(j \cdot d)} \cdot w_{j+1} \right) + b$$This operation is repeated across the entire sequence, producing a resulting vector known as a feature map, $M$. This map represents the degree to which the specific temporal pattern defined by the random kernel is present at each point in the time series. By applying a large ensemble of these kernels (typically 10,000 in the default configuration), ROCKET projects the original time series into a highly diverse, transformed spatial representation.

\subsection{Feature Extraction}
Passing the time series through 10,000 kernels results in 10,000 distinct feature maps. To make this high-dimensional output usable for a standard classifier without causing a combinatorial explosion in memory and processing time, ROCKET applies a pooling step to extract fixed-length scalar summaries from each feature map.For every feature map $M$, ROCKET extracts exactly two aggregate features:Maximum Value (Max): This captures the peak activation of the kernel across the entire time series. It acts as an indicator of the single strongest match between the kernel's pattern and the input data.$$Max(M) = \max_i(M_i)$$Proportion of Positive Values (PPV): This represents the frequency with which a pattern appears in the time series. It calculates the fraction of the feature map where the activation is strictly greater than zero.$$PPV(M) = \frac{1}{|M|} \sum_{i=1}^{|M|} [M_i > 0]$$(where $[ \cdot ]$ represents the Iverson bracket, equating to 1 if the condition is true and 0 otherwise).By extracting two features per kernel, an ensemble of 10,000 kernels produces a final, flat feature vector of 20,000 dimensions for each time series. This extraction mechanism is critical, as it allows ROCKET to handle input series of varying lengths while still outputting a uniform feature vector size for the classifier.

\subsection{Classifier}
Once the feature extraction phase is complete, the time series classification problem is effectively reduced to a standard tabular machine learning task. The 20,000-dimensional feature vectors are used to train a linear classifier.
ROCKET typically employs a Ridge Regression classifier (for smaller datasets) or Logistic Regression (for larger datasets). The choice of a linear model is highly intentional: linear classifiers scale exceptionally well to high-dimensional data and are extremely fast to train. Because the random convolutional kernels combined with the non-linear PPV feature extraction inherently map the data into a space where the classes are largely linearly separable, complex non-linear classifiers (like deep neural networks or complex tree ensembles) are unnecessary. This synergy between diverse, random feature generation and a fast linear backend is the primary driver of ROCKET's computational efficiency.

\subsection{Advantages}
The ROCKET algorithm introduces several distinct advantages over traditional time series classification methods:

\begin{itemize}
    \item Exceptional Computational Efficiency: By eliminating the need to update kernel weights via gradient descent, ROCKET dramatically reduces training time. It achieves state-of-the-art accuracy in a fraction of the time required by deep learning models like InceptionTime or ensemble methods like HIVE-COTE.
    \item High Accuracy: Despite its reliance on random initialization, the sheer volume and diversity of the kernels allow ROCKET to capture a highly discriminative feature set, consistently matching or outperforming computationally heavier benchmarks.
    \item Scalability: The linear nature of the classification phase and the independence of the kernel convolutions make ROCKET highly scalable to both long time series and datasets with massive sample sizes.
    \item Robustness to Hyperparameters: ROCKET requires almost no hyperparameter tuning. The default setting of 10,000 kernels is universally effective across diverse domains, removing the need for costly cross-validation searches.
\end{itemize}

\subsection{Limitations}
While highly effective, the original ROCKET algorithm does have specific constraints:

\begin{itemize}
    \item Memory Consumption: Although computationally fast, storing the 20,000-dimensional feature vectors for every sample in a large dataset can lead to substantial RAM usage, which can be a bottleneck on memory-constrained hardware.
    \item Lack of Interpretability: Due to the random nature of the convolutional kernels, ROCKET acts largely as a "black box." It is exceedingly difficult to trace back exactly which temporal patterns or physical phenomena in the raw data drove the model's final classification decision.
    \item Non-Determinism: The random sampling of weights, biases, and dilations means that feature spaces will vary slightly between runs unless a strict random seed is applied, which can complicate exact reproducibility in some deployment environments.
    \item Univariate Constraint: The base ROCKET algorithm is designed explicitly for univariate time series, requiring architectural modifications (such as those seen in MultiROCKET) to effectively process multivariate data without losing cross-channel context.
\end{itemize}
